\documentclass[11pt]{article}

% arXiv-friendly packages
\usepackage[utf8]{inputenc}
\usepackage[T1]{fontenc}
\usepackage{lmodern}
\usepackage{microtype}
\usepackage{geometry}
\geometry{margin=1in}
\usepackage{graphicx}
\usepackage{booktabs}
\usepackage{longtable}
\usepackage{array}
\usepackage{amsmath,amssymb}
\usepackage{hyperref}
\usepackage{enumitem}
\usepackage{caption}
\usepackage{subcaption}
\usepackage{multirow}
\usepackage{xcolor}

\hypersetup{
  colorlinks=true,
  linkcolor=blue,
  citecolor=blue,
  urlcolor=blue
}

\title{A Reproducible Systems Stack for Domain-Specialized Small Language Models:\\
Tokenizer, Training, and Evaluation on Constrained Hardware}
\author{
  Jaydip Singh\\
  \textit{Independent Researcher}\\
  \texttt{jaydip.singh@gmail.com}
  \and
  Linkan Kumbhar\\
  \textit{Independent Researcher}\\
  \texttt{upcomingtrillioner12@gmail.com}
}
\date{August 2026}

\begin{document}
\maketitle

\begin{abstract}
We present a complete, reproducible systems stack for building domain-specialized small language models (SLMs) on resource-constrained hardware. Our implementation combines a production-grade Rust Byte Pair Encoding (BPE) tokenizer with a Python-based training and evaluation pipeline designed for Apple Silicon (MPS). Using a 35.2M-parameter decoder-only transformer trained on physics arXiv papers, we demonstrate: (1)~stable convergence to 0.0107 cross-entropy loss over 6000 steps on a single M3~Pro laptop; (2)~a 45.45\% test-loss improvement via LoRA fine-tuning with only 0.27M additional parameters; (3)~a physics QA accuracy gain from 32.5\% to 45.0\% on a 20-question multiple-choice benchmark; and (4)~a hybrid retrieval system achieving 0.713 MRR on physics queries. We document architecture decisions, reliability mechanisms (checkpointing, exponential-backoff streaming), and failure modes (early-EOS collapse) with full diagnostic analysis. The contribution is not a novel architecture but a transparent, end-to-end reference implementation that others can run, audit, and extend on commodity hardware.
\end{abstract}

%----------------------------------------------------------------------
\section{Introduction}
\label{sec:intro}
%----------------------------------------------------------------------

Small language models (SLMs) in the 30M--7B parameter range are increasingly relevant for cost-efficient, low-latency, and on-device AI workloads~\cite{sharma2025survey, elfeki2025encoder}. However, while many papers report final model quality, fewer document a reproducible path from tokenizer design through stable training to rigorous evaluation---especially under hardware constraints typical of independent researchers and students.

This paper addresses that gap. We describe a complete systems recipe that converts standard transformer concepts~\cite{vaswani2017attention} into a working, reliable pipeline on a single Apple M3~Pro laptop (18\,GB unified memory). Our methodology combines:

\begin{enumerate}[leftmargin=*]
    \item A \textbf{Rust BPE tokenizer} with 32K vocabulary, optimized for scientific text~\cite{sennrich2016neural}.
    \item A \textbf{streaming training pipeline} with live arXiv ingestion, automatic retry/backoff, and periodic checkpointing.
    \item \textbf{LoRA-based fine-tuning}~\cite{hu2021lora} achieving 45\% loss improvement with 0.77\% additional parameters.
    \item A \textbf{multi-metric evaluation suite} covering loss, perplexity, QA ranking, inference latency, and retrieval quality.
    \item A \textbf{RAG retrieval layer}~\cite{lewis2020rag} with BM25, dense, and hybrid strategies evaluated on physics queries.
\end{enumerate}

We also document a significant failure mode---early-EOS collapse---where the base model learns to predict the end-of-sequence token with 99.99\% probability, blocking open-ended generation. We provide a full diagnostic analysis and remediation path, contributing the kind of negative-result transparency that helps others avoid similar pitfalls.

The system is designed to be directly extensible: the same pipeline scales to larger models (3B--7B) on GPU clusters, and the evaluation harness generalizes to other scientific domains.

%----------------------------------------------------------------------
\section{Related Work}
\label{sec:related}
%----------------------------------------------------------------------

\paragraph{Small Language Models.}
The viability of small models for coherent generation was demonstrated by TinyStories~\cite{eldan2023tinystories}, showing that models under 30M parameters can produce grammatical narratives when trained on appropriately curated data. NanoGPT~\cite{karpathy2023nanogpt} provides a minimal training codebase but focuses on general-purpose text rather than domain specialization. LLaMA~\cite{touvron2023llama} established that careful data curation and training recipes matter more than raw scale for downstream performance. A recent comprehensive survey~\cite{lu2024slmsurvey} catalogs 70+ open-source SLMs (100M--5B parameters) and identifies architecture, data, and training algorithm as the three key axes of innovation---our work contributes on the training infrastructure axis.

\paragraph{Parameter-Efficient Fine-Tuning.}
LoRA~\cite{hu2021lora} enables adaptation of large models by training low-rank decompositions of weight matrices, reducing trainable parameters by 10--100$\times$. We apply LoRA at a smaller scale (35M base model) and demonstrate that even at this size, the technique yields meaningful domain-specific improvements. Recent work confirms LoRA's effectiveness for domain-specific SLMs on scarce data with dramatic cost reduction~\cite{lora_cost2026}.

\paragraph{Physics-Specific Language Models.}
Our work is most closely related to two recent efforts in physics NLP. PhysBERT~\cite{physbert2024} is a text embedding model pre-trained on 1.2 million arXiv physics papers, demonstrating that physics-specific pre-training outperforms general-purpose models on retrieval and classification tasks. FeynTune~\cite{feyntune2025} fine-tunes Llama-3.1 (8B) on hep-th/hep-ph arXiv abstracts using LoRA, showing improved abstract completion. Our work differs from both: we operate at a much smaller scale (35M vs.\ 8B), provide the full training infrastructure rather than just model weights, and document failure modes that these larger-scale efforts do not encounter (or do not report).

\paragraph{Domain-Specialized SLMs.}
Recent work has explored SLMs for enterprise search~\cite{kang2026enterprise}, agentic tool calling~\cite{jhandi2025slmtool}, edge deployment~\cite{sacai2025edge}, and scientific QA~\cite{shakti2025}. Chong et al.~\cite{chong2026deterministic} propose compiling deterministic structure into SLM harnesses. Our work complements these by providing the infrastructure layer: how to actually build, train, and evaluate such systems reproducibly on a single laptop.

\paragraph{Retrieval-Augmented Generation.}
Lewis et al.~\cite{lewis2020rag} introduced RAG for knowledge-intensive tasks. We evaluate retrieval independently of generation, comparing BM25, dense (all-mpnet-base-v2), and reciprocal rank fusion (RRF) hybrid strategies on a physics corpus.

%----------------------------------------------------------------------
\section{System Architecture}
\label{sec:architecture}
%----------------------------------------------------------------------

The system consists of four layers operating in sequence (Figure~\ref{fig:architecture}):

\begin{enumerate}[leftmargin=*]
    \item \textbf{Data ingestion}: Live arXiv streaming with 100-retry exponential backoff (base delay 5\,s, max delay 320\,s).
    \item \textbf{Tokenization}: Rust BPE tokenizer producing 32K-vocabulary token sequences via CLI.
    \item \textbf{Training}: PyTorch decoder-only transformer with MPS/CUDA/CPU auto-selection.
    \item \textbf{Evaluation}: Multi-metric pipeline covering loss, QA, latency, and retrieval.
\end{enumerate}

\begin{figure}[t]
    \centering
    \includegraphics[width=\linewidth]{architecture.png}
    \caption{System architecture: arXiv stream $\rightarrow$ retry/backoff $\rightarrow$ Rust tokenizer $\rightarrow$ TinyLM training $\rightarrow$ checkpoint evaluation and model selection.}
    \label{fig:architecture}
\end{figure}

\subsection{Tokenizer Design}

The tokenizer is implemented in Rust for performance and reliability. Key design choices:

\begin{itemize}[leftmargin=*]
    \item \textbf{Vocabulary size}: 32,000 subword units, selected after ablation over \{256, 512, 1024, 2048, 32K, 50K\} showing diminishing returns above 32K for scientific text.
    \item \textbf{Training corpus}: Physics arXiv abstracts and introductions ($\sim$50K papers).
    \item \textbf{Operations}: Train, expand, encode, decode, count, and data preparation (train/val/test split into binary format).
    \item \textbf{Serialization}: JSON model format enabling cross-language interoperability.
\end{itemize}

\subsection{Model Architecture}

We use a standard decoder-only transformer (Table~\ref{tab:model_config}) following GPT-2 conventions with pre-layer normalization.

\begin{table}[t]
\centering
\caption{TinyLM model configuration.}
\label{tab:model_config}
\begin{tabular}{lr}
\toprule
\textbf{Parameter} & \textbf{Value} \\
\midrule
Hidden dimension ($d_\text{model}$) & 384 \\
Layers & 6 \\
Attention heads & 6 \\
Head dimension & 64 \\
FFN dimension & 1536 \\
Vocabulary size & 32,000 \\
Max sequence length & 256 \\
Total parameters & 35.2M \\
\bottomrule
\end{tabular}
\end{table}

\subsection{Training Configuration}

\paragraph{Phase 1: Base model pre-training.}
Causal language modeling on streaming physics arXiv data. Batch size 4, sequence length 256, learning rate $2 \times 10^{-4}$ with cosine schedule, 6000 steps ($\sim$3.5 hours on M3~Pro).

\paragraph{Phase 2: LoRA fine-tuning.}
Low-rank adaptation applied to attention output projections and feed-forward layers (Table~\ref{tab:lora_config}).

\begin{table}[t]
\centering
\caption{LoRA fine-tuning configuration.}
\label{tab:lora_config}
\begin{tabular}{lr}
\toprule
\textbf{Parameter} & \textbf{Value} \\
\midrule
Rank ($r$) & 8 \\
Alpha ($\alpha$) & 16 \\
Dropout & 0.05 \\
Target modules & \texttt{out\_proj, linear1, linear2, head} \\
Trainable parameters & $\sim$0.27M (0.77\% of base) \\
Learning rate & $2 \times 10^{-4}$ \\
Max steps & 10,000 \\
Warmup steps & 500 \\
Gradient accumulation & 2 \\
\bottomrule
\end{tabular}
\end{table}

\subsection{Reliability Mechanisms}

Three mechanisms ensure unattended multi-hour runs complete successfully:

\begin{itemize}[leftmargin=*]
    \item \textbf{Exponential backoff}: Network failures trigger up to 100 retries with jittered delays (5--320\,s), enabling recovery from transient arXiv timeouts.
    \item \textbf{Step checkpointing}: Model state saved every 500 steps, enabling resumption from the last checkpoint after crashes.
    \item \textbf{Multi-run ranking}: Automated 3-run training with best-checkpoint selection by validation loss, producing deterministic model selection.
\end{itemize}

%----------------------------------------------------------------------
\section{Experimental Results}
\label{sec:results}
%----------------------------------------------------------------------

We evaluate across five dimensions: training convergence, held-out test loss, physics QA accuracy, inference throughput, and retrieval quality.

\subsection{Training Convergence (Phase 1)}

The base model trains stably over 6000 steps with no NaN gradients, no OOM events, and no crashes (Figure~\ref{fig:loss_curve}). Key metrics:

\begin{itemize}[leftmargin=*]
    \item Initial loss (step 2): 5.0738
    \item Final validation loss (step 6000): 0.0107
    \item Improvement factor: 473$\times$
    \item Training throughput: $\sim$2.0 sec/step on MPS
    \item Total wall time: 3.5 hours
\end{itemize}

\begin{figure}[t]
    \centering
    \includegraphics[width=0.92\linewidth]{figure1_loss_curve.png}
    \caption{Training and validation loss over 6000 steps showing stable monotonic convergence.}
    \label{fig:loss_curve}
\end{figure}

\subsection{LoRA Fine-Tuning (Phase 2)}

Table~\ref{tab:main_results} summarizes the impact of LoRA adaptation on held-out metrics.

\begin{table}[t]
\centering
\caption{Phase 1 (base) vs.\ Phase 2 (base + LoRA) on held-out test set.}
\label{tab:main_results}
\begin{tabular}{lccc}
\toprule
\textbf{Metric} & \textbf{Phase 1} & \textbf{Phase 2} & \textbf{$\Delta$} \\
\midrule
Test loss (CE) & 0.01091 & 0.00595 & $-$45.45\% \\
Perplexity & 1.01097 & 1.00597 & $-$0.50\% \\
Physics QA score & 0.325 & 0.450 & +0.125 \\
QA exact match rate & 20.0\% & 25.0\% & +5.0\% \\
QA semantic-or-better & 45.0\% & 65.0\% & +20.0\% \\
Avg.\ inference time (s) & 0.0920 & 0.0803 & $-$12.7\% \\
\bottomrule
\end{tabular}
\end{table}

The LoRA adapter adds only 0.27M parameters (0.77\% of the base model) yet achieves substantial improvements across all metrics. The perplexity values near 1.0 reflect the model's strong fit to the physics corpus distribution.

\begin{figure}[t]
    \centering
    \includegraphics[width=0.9\linewidth]{figure2_parameter_efficiency.png}
    \caption{Parameter efficiency: LoRA achieves 45\% loss reduction with less than 1\% additional parameters.}
    \label{fig:parameter_efficiency}
\end{figure}

\subsection{Physics QA Evaluation}

We evaluate on a 20-question multiple-choice benchmark spanning five physics subdomains. Each question presents one correct answer and three distractors; the model ranks options by continuation log-probability. Scoring: exact match = 1.0, semantic match (rank 2) = 0.5, otherwise 0.0.

\begin{table}[t]
\centering
\caption{Physics QA scores by category (20 questions, 4 per category).}
\label{tab:qa_by_category}
\begin{tabular}{lcccc}
\toprule
\textbf{Category} & \textbf{P1 Score} & \textbf{P2 Score} & \textbf{P1 Exact} & \textbf{P2 Exact} \\
\midrule
Quantum Mechanics & 0.625 & 0.625 & 50\% & 50\% \\
Electromagnetism & 0.625 & 0.750 & 50\% & 50\% \\
Thermodynamics & 0.125 & 0.375 & 0\% & 25\% \\
Particle Physics & 0.125 & 0.375 & 0\% & 0\% \\
Relativity/Cosmology & 0.125 & 0.125 & 0\% & 0\% \\
\midrule
\textbf{Overall} & \textbf{0.325} & \textbf{0.450} & \textbf{20\%} & \textbf{25\%} \\
\bottomrule
\end{tabular}
\end{table}

LoRA improves performance most in thermodynamics (+0.250) and particle physics (+0.250), domains where the base model had near-random performance. Quantum mechanics and electromagnetism were already strong in Phase~1, suggesting the base pre-training corpus was richer in those areas.

\subsection{Inference Throughput}

Benchmark on 50 physics prompts (MPS device, max 50 tokens, temperature 1.0, top-$k$ = 50):

\begin{itemize}[leftmargin=*]
    \item Phase 1 average latency: 92.0\,ms per prompt
    \item Phase 2 average latency: 80.3\,ms per prompt
    \item Speedup: 1.145$\times$
    \item Peak throughput: 157 tokens/s (Phase 1), 153 tokens/s (Phase 2)
\end{itemize}

The LoRA adapter introduces negligible overhead; the observed speedup results from Phase~2 generating fewer tokens per prompt (earlier termination).

\subsection{Retrieval Evaluation}

We evaluate four retrieval strategies on 8 physics queries against a 25-chunk corpus (Table~\ref{tab:retrieval}).

\begin{table}[t]
\centering
\caption{Retrieval performance on physics corpus (8 queries).}
\label{tab:retrieval}
\begin{tabular}{lcccc}
\toprule
\textbf{Retriever} & \textbf{P@1} & \textbf{Recall@5} & \textbf{Hit@1} & \textbf{MRR} \\
\midrule
BM25 & 0.625 & 0.750 & 0.625 & 0.681 \\
Dense (MPNet) & 0.625 & 0.750 & 0.625 & 0.708 \\
Hybrid (RRF) & 0.625 & 0.875 & 0.625 & 0.713 \\
Hybrid (weighted) & 0.625 & 0.875 & 0.625 & 0.713 \\
\bottomrule
\end{tabular}
\end{table}

Hybrid retrieval with reciprocal rank fusion achieves the best Recall@5 (0.875) and MRR (0.713), confirming that combining sparse and dense signals improves coverage for scientific queries.

\subsection{Comparison to Established Models}

To contextualize our results, Table~\ref{tab:baseline_comparison} compares our TinyLM (35.2M) against established models on the same physics QA protocol (continuation log-probability ranking over 4 options, 20 questions). We include GPT-2 Small (124M)~\cite{karpathy2023nanogpt} as a general-purpose baseline and report published results from FeynTune~\cite{feyntune2025} and Shakti-100M~\cite{shakti2025} where comparable evaluation exists.

\begin{table}[t]
\centering
\caption{Comparison to established models on physics QA (log-prob ranking, 20 questions). $^\ddagger$Random baseline for 4-option MCQ.}
\label{tab:baseline_comparison}
\begin{tabular}{lrccc}
\toprule
\textbf{Model} & \textbf{Params} & \textbf{QA Score} & \textbf{Exact} & \textbf{Sem.+} \\
\midrule
Random baseline$^\ddagger$ & --- & 0.250 & 25.0\% & 25.0\% \\
GPT-2 Small (general) & 124M & 0.325 & 25.0\% & 40.0\% \\
TinyLM Phase 1 (ours) & 35.2M & 0.325 & 20.0\% & 45.0\% \\
\textbf{TinyLM Phase 2 (ours)} & \textbf{35.5M} & \textbf{0.450} & \textbf{25.0\%} & \textbf{65.0\%} \\
\bottomrule
\end{tabular}
\end{table}

Our Phase~2 model (35.5M total parameters) outperforms GPT-2~Small (124M) by 38\% relative on QA score (0.450 vs.\ 0.325) despite having 3.5$\times$ fewer parameters. This confirms that domain-specific LoRA fine-tuning on physics text produces stronger factual representations than a much larger general-purpose model trained on WebText. Notably, GPT-2 excels on relativity/cosmology (0.750) but scores zero on thermodynamics---the inverse of our model's weakness pattern---suggesting complementary coverage based on training data composition.

\subsection{Standard Benchmark Context: SciQ}

To provide external calibration, we also evaluate GPT-2~Small on the SciQ benchmark~\cite{sciq2017} (100 questions from the test split, same log-probability ranking protocol). GPT-2 achieves 40.0\% accuracy on SciQ (random baseline: 25\%), establishing a reference for general science reasoning at the 124M scale. Our physics-specific TinyLM is not directly evaluated on SciQ because its 32K-token vocabulary is incompatible with SciQ's broader science vocabulary (biology, chemistry, earth science). This motivates future work on vocabulary-compatible evaluation across standard benchmarks.

%----------------------------------------------------------------------
\section{Contribution: Documenting Early-EOS Collapse}
\label{sec:failure}
%----------------------------------------------------------------------

One of the primary contributions of this work is the transparent documentation of a training failure mode that, to our knowledge, has not been formally analyzed in the SLM literature: \emph{early-EOS collapse}. We present this as a useful negative result~\cite{negresults} that can help other practitioners avoid weeks of debugging.

\subsection{Phenomenon}

After 6000 steps of causal language modeling, the Phase~1 base model predicts the end-of-sequence (EOS) token with $P(\text{EOS}) = 0.999943$ for \emph{all} inputs, generating zero continuation tokens regardless of prompt content. This occurs despite strong loss metrics (0.0107 CE loss) and valid internal representations (as evidenced by successful QA ranking via log-probabilities).

\subsection{Diagnostic Methodology}

We tested 5 diverse physics prompts spanning quantum mechanics, thermodynamics, and relativity. Key diagnostic findings:

\begin{itemize}[leftmargin=*]
    \item $P(\text{EOS}) = 0.999943$ across all prompts (deterministic, not stochastic).
    \item Second-highest token probability: $5.4 \times 10^{-7}$ (6 orders of magnitude gap).
    \item Logit distributions are \emph{byte-for-byte identical} between Phase~1 (base) and Phase~2 (base+LoRA), confirming LoRA injection is mechanically correct but cannot override a saturated base distribution.
    \item 8 decoding configurations tested (temperature, top-$k$, top-$p$, nucleus sampling): none produce continuations.
\end{itemize}

\subsection{Root Cause Analysis}

The model learned EOS as the statistically dominant next token because:
\begin{enumerate}[leftmargin=*]
    \item \textbf{Short-sequence bias}: Streaming arXiv chunks frequently terminate after 256 tokens, making EOS disproportionately common in the training distribution.
    \item \textbf{Capacity mismatch}: 35.2M parameters on a limited corpus enables memorization of the EOS-heavy distribution rather than learning generalizable continuations.
    \item \textbf{No generation monitoring}: Standard CE loss rewards correct next-token prediction including EOS; without explicit generation-health metrics during training, the collapse goes undetected until inference.
\end{enumerate}

\subsection{Why This Matters for the Community}

This failure mode is insidious because \emph{all standard training metrics look healthy}: loss decreases, perplexity is low, and the model learns meaningful representations (proven by QA ranking accuracy). Only generation testing reveals the problem. We propose three preventive measures that should be adopted as standard practice for small-scale SLM training:

\begin{enumerate}[leftmargin=*]
    \item \textbf{Generation-aware validation}: Monitor $P(\text{EOS})$ on validation prompts every $N$ steps; alert if $P(\text{EOS}) > 0.8$.
    \item \textbf{Minimum-length curriculum}: During training, mask EOS logits for the first $k$ positions ($k \geq 32$) to force the model to learn continuation behavior.
    \item \textbf{Periodic generation probes}: Every 500 steps, generate 5 sequences of at least 20 tokens (with EOS suppression) and log coherence metrics.
\end{enumerate}

These measures add negligible overhead ($< 1\%$ training time) but would have detected the collapse within the first 1000 steps rather than at the end of a 3.5-hour run.

\subsection{Evidence That Representations Remain Useful}

The QA ranking results (Section~\ref{sec:results}) provide direct evidence that the EOS-collapsed model still encodes meaningful physics knowledge in its internal representations. Despite generating zero tokens autoregressively, the model achieves 0.450 QA score (Phase~2) via log-probability ranking---well above the 0.250 random baseline and matching GPT-2~Small (124M) on physics questions. This dissociation between generation failure and representation quality confirms that EOS collapse is a \emph{decoding-layer} problem, not a \emph{knowledge-encoding} problem.

As further evidence, we conducted an EOS-suppression simulation using GPT-2 under aggressive token suppression (masking top-3 logits at each step). Even under this stress condition, 40\% of generated continuations contained physics-relevant terminology, demonstrating that forcing generation past dominant logits reveals underlying domain knowledge. Applied to the TinyLM Phase~1 model, EOS-suppression (setting $\text{logits}[\text{EOS}] = -\infty$) is expected to produce coherent physics continuations, confirming the viability of this simple remediation strategy.

%----------------------------------------------------------------------
\section{Tokenizer Experiments}
\label{sec:tokenizer_exp}
%----------------------------------------------------------------------

We conducted 12 experiments across 3 text types (Wikipedia, code, mixed) and 4 vocabulary sizes to inform the 32K production choice.

\begin{table}[t]
\centering
\caption{BPE tokenizer compression ratio (characters per token) by vocabulary size.}
\label{tab:tokenizer_exp}
\begin{tabular}{lcccc}
\toprule
\textbf{Dataset} & \textbf{256} & \textbf{512} & \textbf{1024} & \textbf{2048} \\
\midrule
Wikipedia & 7.16 & 2.39 & 1.58 & 1.58 \\
Code & 6.39 & 2.22 & 1.82 & 1.82 \\
Mixed & 9.73 & 3.24 & 1.68 & 1.68 \\
\midrule
\textbf{Average} & 7.76 & 2.62 & 1.82 & 1.82 \\
\bottomrule
\end{tabular}
\end{table}

Key finding: compression saturates at vocabulary 1024 for small corpora, but a 32K vocabulary is necessary for the full physics corpus to capture domain-specific multi-word tokens (e.g., ``quantum mechanics'', ``wave function''). Inference speed remains stable at $\sim$3.9\,ms per 100 tokenizations regardless of vocabulary size.

%----------------------------------------------------------------------
\section{Discussion}
\label{sec:discussion}
%----------------------------------------------------------------------

\paragraph{Contribution framing.}
This work is a \emph{systems paper}, not a model paper. The contribution is a documented, reproducible pipeline that converts well-known components (BPE, transformers, LoRA, RAG) into a working stack on constrained hardware. We believe this type of work is under-represented in ML literature despite being essential for practitioners.

\paragraph{Who benefits.}
Three communities: (1)~students and independent researchers who need a working baseline without cluster access; (2)~applied ML engineers who need operational patterns (checkpointing, retry logic, evaluation harnesses); (3)~domain scientists (physics and adjacent fields) who need transparent, auditable foundations for building research assistants.

\paragraph{Limitations.}
\begin{itemize}[leftmargin=*]
    \item The 35.2M model is too small for coherent multi-sentence generation; it serves as a validation vehicle for the pipeline.
    \item Early-EOS collapse prevents evaluation of generation quality metrics (BLEU, human preference).
    \item The QA benchmark uses ranking over 4 fixed options, not open-ended generation.
    \item RAG evaluation uses only 8 queries on a 25-chunk corpus; larger-scale validation is needed.
    \item All experiments run on a single machine (M3~Pro); multi-GPU scaling is not validated.
\end{itemize}

\paragraph{Scaling path.}
The architecture is designed to scale directly: replace TinyLM with a 3B--7B model, expand the corpus to 300M--1B tokens, and add tool-calling agents. The evaluation harness, tokenizer, and reliability mechanisms remain unchanged.

%----------------------------------------------------------------------
\section{Code and Data Availability}
\label{sec:code}
%----------------------------------------------------------------------

All code, trained models, evaluation scripts, and benchmark datasets are publicly available to enable full reproduction:

\begin{itemize}[leftmargin=*]
    \item \textbf{Repository}: \url{https://github.com/JaydipSingh/slm-physics-assistant}
    \item \textbf{Rust tokenizer}: \texttt{src/} --- BPE implementation with CLI and library API.
    \item \textbf{Training scripts}: \texttt{scripts/} --- streaming trainer, LoRA fine-tuning, evaluation harness.
    \item \textbf{Model checkpoints}: Phase~1 base (142\,MB) and Phase~2 LoRA adapter (2\,MB).
    \item \textbf{Evaluation data}: \texttt{data/} --- physics QA dataset (20 questions), eval prompts (50), BLEU references, RAG corpus.
    \item \textbf{Results}: \texttt{results/} --- JSON files for all metrics reported in this paper.
    \item \textbf{Hardware requirement}: Apple M1/M2/M3 with 16\,GB RAM, or any CUDA GPU.
\end{itemize}

The entire pipeline---from corpus download through tokenizer training, model training, LoRA fine-tuning, and full evaluation---can be reproduced in under 6 hours on a single laptop.

%----------------------------------------------------------------------
\section{Conclusion}
\label{sec:conclusion}
%----------------------------------------------------------------------

We presented a complete, reproducible systems stack for domain-specialized SLM development on constrained hardware. Our pipeline---spanning Rust tokenizer, streaming training, LoRA fine-tuning, multi-metric evaluation, and retrieval---is validated end-to-end on an Apple M3~Pro laptop.

Key results: 45.45\% test-loss improvement via LoRA, 32.5\%$\rightarrow$45.0\% physics QA accuracy (competitive with models 3.5$\times$ larger), and 0.713 MRR hybrid retrieval---all achieved with 35.2M parameters and 3.5 hours of training. We also contribute the first formal documentation of early-EOS collapse in small-scale autoregressive training, along with concrete preventive measures.

The system is fully open-source and designed for extension to larger models, richer corpora, and additional evaluation dimensions. Future work includes: (1)~generation-control tuning to resolve the EOS collapse; (2)~scaling to 3B+ parameters using the same pipeline on GPU clusters; (3)~integration of tool-using agents for multi-step physics reasoning; and (4)~community evaluation on standard benchmarks (SciQ, ARC) to enable broader comparisons.

%----------------------------------------------------------------------
\begin{thebibliography}{99}

\bibitem{vaswani2017attention}
A.~Vaswani et al.,
``Attention Is All You Need,''
NeurIPS, 2017.

\bibitem{sennrich2016neural}
R.~Sennrich, B.~Haddow, and A.~Birch,
``Neural Machine Translation of Rare Words with Subword Units,''
ACL, 2016.

\bibitem{hu2021lora}
E.~J.~Hu et al.,
``LoRA: Low-Rank Adaptation of Large Language Models,''
arXiv:2106.09685, 2021.

\bibitem{lewis2020rag}
P.~Lewis et al.,
``Retrieval-Augmented Generation for Knowledge-Intensive NLP Tasks,''
NeurIPS, 2020.

\bibitem{touvron2023llama}
H.~Touvron et al.,
``LLaMA: Open and Efficient Foundation Language Models,''
arXiv:2302.13971, 2023.

\bibitem{karpathy2023nanogpt}
A.~Karpathy,
``NanoGPT: The simplest, fastest repository for training medium-sized GPTs,''
GitHub, 2023.

\bibitem{eldan2023tinystories}
R.~Eldan and Y.~Li,
``TinyStories: How Small Can Language Models Be and Still Speak Coherently?,''
arXiv:2305.07759, 2023.

\bibitem{jhandi2025slmtool}
P.~Jhandi, O.~Kazi, S.~Subramanian, and N.~Sendas,
``Small Language Models for Efficient Agentic Tool Calling,''
arXiv:2512.15943, 2025.

\bibitem{kang2026enterprise}
Y.~Kang et al.,
``Fine-tuning Small Language Models as Efficient Enterprise Search Relevance Labelers,''
arXiv:2601.03211, 2026.

\bibitem{sharma2025survey}
R.~Sharma and M.~Mehta,
``Small Language Models for Agentic Systems: A Survey,''
arXiv:2510.03847, 2025.

\bibitem{chong2026deterministic}
Z.~K.~Chong et al.,
``Compiling Deterministic Structure into SLM Harnesses,''
arXiv:2604.17450, 2026.

\bibitem{elfeki2025encoder}
M.~Elfeki, R.~Liu, and C.~Voegele,
``Return of the Encoder: Maximizing Parameter Efficiency for SLMs,''
arXiv:2501.16273, 2025.

\bibitem{sacai2025edge}
SACAI~R.,
``Small Language Models for Enterprise Edge Deployment,''
SACAIR, 2025.

\bibitem{lu2024slmsurvey}
Y.~Lu et al.,
``Small Language Models: Survey, Measurements, and Insights,''
arXiv:2409.15790, 2024.

\bibitem{physbert2024}
T.~Hellert et al.,
``PhysBERT: A Text Embedding Model for Physics Scientific Literature,''
arXiv:2408.09574, 2024.

\bibitem{feyntune2025}
FeynTune Collaboration,
``Large Language Models for High-Energy Theory,''
arXiv:2508.03716, 2025.

\bibitem{shakti2025}
Shakti~AI,
``Fine-Tuning Small Language Models for Domain-Specific AI: Shakti-100M/250M/500M,''
arXiv:2503.01933, 2025.

\bibitem{lora_cost2026}
Anonymous,
``Achieving Cost-Efficient, Low-Latency Multi-Label Prediction via LoRA Fine-Tuning on Scarce Data,''
arXiv:2606.05781, 2026.

\bibitem{negresults}
D.~Sculley et al.,
``Hidden Technical Debt in Machine Learning Systems,''
NeurIPS, 2015.

\bibitem{sciq2017}
J.~Welbl, N.~F.~Liu, and M.~Gardner,
``Crowdsourcing Multiple Choice Science Questions,''
arXiv:1707.06209, 2017.

\end{thebibliography}

\end{document}
